\section{I. INTRODUCTION}

\subsection{Paragraf 1: Landasan Kontekstual}
\begin{frame}{Paragraf 1 (1-6)}
    \begin{itemize}
        \item \textbf{1. finance $\to$ derivative:} Mendefinisikan instrumen derivatif secara luas sebagai fondasi kontekstual.
        \item \textbf{2. jenis khusus:} Spesifikasi pada derivatif suku bunga sebagai fokus utama.
        \item \textbf{3. tujuan derivatif:} Utilitas praktis dalam manajemen risiko (\textit{hedging}) dan spekulasi.
        \item \textbf{4. pemodelan matematis:} Akurasi model evolusi suku bunga sebagai prasyarat \textit{pricing}.
        \item \textbf{5. jenis model:} Klasifikasi antara \textit{spot rate} dan \textit{forward rate}.
        \item \textbf{6. model simpel:} Latar belakang historis penggunaan faktor \textit{noisy} tunggal atau ganda.
    \end{itemize}
\end{frame}

\begin{frame}{Paragraf 1 (7-12)}
    \begin{itemize}
        \item \textbf{7. spesifik model:} Identifikasi model klasik (Vasicek, Hull-White, CIR) sebagai standar industri.
        \item \textbf{8. kekurangan model (1):} Mengakui kemudahan implementasi namun beralih ke limitasi performa.
        \item \textbf{9. kekurangan model (2):} Kritik pada kegagalan kalibrasi data pasar dan penangkapan korelasi.
        \item \textbf{10. contoh model bagus:} Pengenalan kerangka kerja \textit{Heath-Jarrow-Morton} (HJM).
        \item \textbf{11. definisi 1:} Otoritas HJM sebagai \textit{general family} model finansial.
        \item \textbf{12. definisi 2:} Referensi ke Appendix A untuk detail teoretis mendalam.
    \end{itemize}
\end{frame}

\begin{frame}{Paragraf 1 (13-17)}
    \begin{itemize}
        \item \textbf{13. definisi 3:} Mereduksi HJM ke faktor volatilitas sebagai target algoritma kuantum.
        \item \textbf{14. kekurangan:} \textit{Trade-off} antara jumlah faktor \textit{noisy} dan efisiensi waktu eksekusi klasik.
        \item \textbf{15. kekurangan:} Limitasi daya komputasi membatasi akurasi dan skalabilitas model.
        \item \textbf{16. QC sebagai solusi:} Komputasi kuantum sebagai paradigma baru penembus hambatan kapasitas.
        \item \textbf{17. lanjut solusi:} Proyeksi keunggulan kuantum untuk fidelitas model HJM.
    \end{itemize}
\end{frame}

\begin{frame}{Pertanyaan yang dijawab pada Paragraf 1}
    \centering
    \small
    \begin{tabular}{lp{4.5cm}p{6cm}}
        \toprule
        Kalimat ke & Pertanyaan & Jawaban \\
        \midrule
        1 & Apa entitas fundamental? & Derivatif sebagai kontrak dari aset dasar. \\
        4 & Prasyarat utama pricing? & Pemodelan akurat evolusi waktu suku bunga. \\
        10 & Kerangka kerja solusi? & \textit{Heath-Jarrow-Morton} (HJM). \\
        13 & Parameter kunci HJM? & Seluruh dinamika ditentukan oleh faktor volatilitas. \\
        16 & Instrumen solusi? & Komputer kuantum (\textit{Quantum Computer}). \\
        \bottomrule
    \end{tabular}
\end{frame}

\subsection{Paragraf 2: Komputasi Kuantum}
\begin{frame}{Paragraf 2}
    \begin{itemize}
        \item \textbf{1. kontribusi QC:} Konteks historis QC sebagai teknologi disruptif revolusioner.
        \item \textbf{2. entanglement:} Sumber daya unik untuk paralelisasi dan potensi \textit{speedup}.
        \item \textbf{3. lanjutan:} Bukti empiris algoritma kuantum (faktorisasi, \textit{searching}, \textit{eigenvalue}).
        \item \textbf{4. quantum simulation:} Aplikasi QS untuk sistem yang inefisien jika disimulasikan klasik.
        \item \textbf{5. conth qs:} Spektrum QS di fisika, kimia, dan biologi.
        \item \textbf{6. paper ccotnoh qs:} Tinjauan literatur domain keuangan dan \textit{research gap} implementasi.
        \item \textbf{7. batasan qs:} Keterbatasan teknologi saat ini (\textit{small noisy chips}) sebagai urgensi riset.
    \end{itemize}
\end{frame}

\begin{frame}{Pertanyaan yang dijawab pada Paragraf 2}
    \centering
    \small
    \begin{tabular}{lp{4.5cm}p{6cm}}
        \toprule
        Kalimat ke & Pertanyaan & Jawaban \\
        \midrule
        2 & Apa peran \textit{entanglement}? & Sumber daya ekstra untuk mempercepat kinerja. \\
        4 & Aplikasi paling relevan? & Simulasi kuantum (QS) untuk sistem inefisien. \\
        6 & Status di finansial? & Diusulkan teoretis, eksperimen riil terbatas. \\
        7 & Hambatan hardware? & Hanya sedia chip kuantum kecil dan berderau. \\
        \bottomrule
    \end{tabular}
\end{frame}

\subsection{Paragraf 3: Implementasi qPCA}
\begin{frame}{Paragraf 3}
    \begin{itemize}
        \item \textbf{1. pengenalan qPCA:} qPCA sebagai kontribusi utama mereduksi kompleksitas model HJM.
        \item \textbf{2. IBMQX2 5-qubit:} Rincian teknis implementasi hardware untuk menunjukkan validitas eksperimental.
        \item \textbf{3. volatilitas \& eigen:} Hubungan parameter fisik dengan representasi \textit{eigenvalue/eigenvector}.
        \item \textbf{6. Quantum Monte Carlo:} qPCA sebagai langkah kritis pembangunan algoritma QMC yang efisien.
        \item \textbf{7. era NISQ:} Modifikasi algoritma agar adaptif terhadap limitasi koherensi prosesor.
    \end{itemize}
\end{frame}

\begin{frame}{Pertanyaan yang dijawab pada Paragraf 3}
    \centering
    \small
    \begin{tabular}{lp{4.5cm}p{6cm}}
        \toprule
        Kalimat ke & Pertanyaan & Jawaban \\
        \midrule
        1 & Solusi artikel ini? & Algoritma qPCA yang efisien mereduksi faktor \textit{noisy}. \\
        3 & Representasi finansial? & Volatilitas diestimasi melalui \textit{eigenvalue/vector}. \\
        6 & Mengapa reduksi krusial? & Meminimalkan jumlah gerbang dalam algoritma QMC. \\
        \bottomrule
    \end{tabular}
\end{frame}

\subsection{Paragraf 4: Definisi PCA dan Speekup}
\begin{frame}{Paragraf 4}
    \begin{itemize}
        \item \textbf{1. definisi PCA:} Teknik matematis mencari aproksimasi matriks \textit{low-rank} optimal.
        \item \textbf{2. reduksi dimensi:} Mekanisme eliminasi \textit{eigenvalue} terkecil untuk mempertahankan informasi krusial.
        \item \textbf{4. biaya komputasi:} Kendala PCA klasik berupa lonjakan biaya pada matriks besar.
        \item \textbf{5. speedup kuantum:} Algoritma kuantum menawarkan \textit{exponential speedup} untuk PCA.
        \item \textbf{7. asumsi teknis:} Matriks direpresentasikan sebagai \textit{quantum state} (\textit{trace} = 1).
    \end{itemize}
\end{frame}

\begin{frame}{Pertanyaan yang dijawab pada Paragraf 4}
    \centering
    \small
    \begin{tabular}{lp{4.5cm}p{6cm}}
        \toprule
        Kalimat ke & Pertanyaan & Jawaban \\
        \midrule
        1 & Definisi teknik PCA? & Aproksimasi matriks \textit{low-rank} yang optimal. \\
        4 & Kendala PCA klasik? & Biaya komputasi tinggi saat ukuran matriks naik. \\
        5 & Peran QC di PCA? & Menawarkan \textit{speedup} eksponensial. \\
        \bottomrule
    \end{tabular}
\end{frame}

\subsection{Paragraf 5: Kecocokan qPCA Finansial}
\begin{frame}{Paragraf 5}
    \begin{itemize}
        \item \textbf{1. eksplorasi qPCA:} Rencana reduksi dimensi model HJM tanpa mengorbankan akurasi.
        \item \textbf{2. density matrices:} Struktur matriks korelasi finansial sangat cocok dengan representasi status kuantum.
        \item \textbf{3. mitigasi dekoherensi:} Reduksi gerbang sirkuit krusial untuk mengatasi limitasi waktu koherensi.
    \end{itemize}
\end{frame}

\begin{frame}{Pertanyaan yang dijawab pada Paragraf 5}
    \centering
    \small
    \begin{tabular}{lp{4.5cm}p{6cm}}
        \toprule
        Kalimat ke & Pertanyaan & Jawaban \\
        \midrule
        2 & Mengapa qPCA cocok? & Karena input identik karakteristik \textit{density matrices}. \\
        3 & Dampak reduksi? & Mengatasi limitasi waktu dekoherensi prosesor. \\
        \bottomrule
    \end{tabular}
\end{frame}

\subsection{Paragraf 6: Visi dan Klaim Kebaruan}
\begin{frame}{Paragraf 6}
    \begin{itemize}
        \item \textbf{1. modifikasi algoritma:} Adaptasi versi standar untuk chip kuantum kecil dan berderau.
        \item \textbf{2. langkah awal:} Fondasi menuju simulasi penuh model HJM pada komputer kuantum IBM.
        \item \textbf{3. klaim novelty:} Eksperimen pertama dalam \textit{option pricing} dan implementasi qPCA terbesar.
        \item \textbf{4. visi strategis:} Membuka jalan bagi pencapaian \textit{quantum supremacy} di bidang keuangan.
    \end{itemize}
\end{frame}

\begin{frame}{Pertanyaan yang dijawab pada Paragraf 6}
    \centering
    \small
    \begin{tabular}{lp{4.5cm}p{6cm}}
        \toprule
        Kalimat ke & Pertanyaan & Jawaban \\
        \midrule
        1 & Karakteristik algoritma? & Modifikasi agar adaptif chip kecil dan berderau. \\
        3 & Klaim kebaruan? & Eksperimen \textit{option pricing} pertama via QC. \\
        4 & Visi jangka panjang? & Pencapaian \textit{quantum advantage} di keuangan. \\
        \bottomrule
    \end{tabular}
\end{frame}
